%% Repairing faulty logical theories in Isabelle/HOL with the operations of
%% the ABC system.
%% The target journal has not been chosen yet. The elsarticle class is used as
%% a neutral preprint format.
\documentclass[preprint,12pt]{elsarticle}

\usepackage[utf8]{inputenc}
\usepackage[T1]{fontenc}
\usepackage{lmodern}
\usepackage{amsmath,amssymb,amsthm}
\usepackage{booktabs,array,tabularx}
\usepackage{xcolor}
\usepackage{listings}
\usepackage{url}
\usepackage[hidelinks]{hyperref}

\newtheorem{definition}{Definition}

\newcommand{\isa}[1]{\texttt{#1}}
\newcommand{\at}[1]{\mathit{#1}}
\newcommand{\TPS}{T(\mathit{PS})}
\newcommand{\FPS}{F(\mathit{PS})}
\newcommand{\imp}{\Longrightarrow}

\definecolor{kw}{RGB}{0,0,160}
\lstdefinelanguage{Isar}{
  morekeywords={theory,imports,begin,end,fun,where,lemma,theorem,by,
    abbreviation,abc_faults,abc_repair,axioms,true,false,ML,val,fn,
    datatype,type,of},
  sensitive=true,
  morecomment=[n]{(*}{*)},
  morestring=[b]",
  literate={⟹}{{$\Longrightarrow$}}2 {∧}{{$\wedge$}}1 {≠}{{$\neq$}}1
    {⇒}{{$\Rightarrow$}}1 {≡}{{$\equiv$}}1 {≤}{{$\le$}}1
    {⟷}{{$\longleftrightarrow$}}2
    {‹}{{$\langle$}}1 {›}{{$\rangle$}}1
}
\lstdefinelanguage{Out}{
  sensitive=true,
  basicstyle=\ttfamily\footnotesize,
  literate={⟹}{{$\Longrightarrow$}}2 {∧}{{$\wedge$}}1 {≠}{{$\neq$}}1
}
\lstset{language=Isar,basicstyle=\ttfamily\small,keywordstyle=\color{kw}\bfseries,
  columns=fullflexible,keepspaces=true,showstringspaces=false,upquote=true,
  frame=single,framesep=3pt,
  xleftmargin=4pt,xrightmargin=4pt,breaklines=true}

\journal{}

\begin{document}

\begin{frontmatter}

\title{Repairing faulty logical theories in Isabelle/HOL with the operations
of the ABC system}

\author[aff]{Omar Monta\~{n}o-Rivas}
% TODO: complete the affiliation and the corresponding author's e-mail.
\address[aff]{Affiliation to be completed}

\begin{abstract}
The ABC system of Li, Bundy and Smaill repairs a logical theory whose
predictions disagree with observations. An observed fact that is not a
theorem reveals an insufficiency, and a theorem that contradicts an
observation reveals an incompatibility. ABC analyses the failed and the
unwanted proofs and repairs the theory with operations drawn from
abduction, belief revision and conceptual change, some of which change the
language of the theory. We describe an implementation of the fault
diagnosis and of the ten repair operations of ABC as an extension of
IsaScheme, a theory exploration system for Isabelle/HOL. Theories are sets
of Horn clauses written in Isabelle's syntax, terms may contain function
symbols, and proofs are searched by selected literal resolution up to a
depth limit. A breadth-first search reports, for each number of
operations, the fault-free theories that the operations produce. The
extension reproduces the repairs of the worked examples of ABC, from the
black swan to the root cause analysis of a 5G network and the diagnosis of
a student's arithmetic misconception. It also shows that, with the complete
table of digit facts, the repair that forgets the carry predicts any tens
digit, which a single answer in the false set exposes. Finally, we carry the
arithmetic example over to recursive functions in higher-order logic. There
the repair removes an argument of a function definition, and IsaScheme's
inductive prover proves what the student's procedure computes and why the
misconception is hard to notice.
\end{abstract}

\begin{keyword}
theory repair \sep abduction \sep belief revision \sep reformation \sep
Horn clauses \sep theory exploration \sep Isabelle/HOL
\end{keyword}

\end{frontmatter}

%% ===========================================================================
\section{Introduction}
\label{sec:intro}

A logical theory that describes an environment makes predictions, which
are its theorems. When the environment is observed, some predictions may
turn out to be false and some observed facts may not be predicted at all.
The ABC system~\cite{LiBundySmaill18,Li21,Bundy26} diagnoses both kinds of
fault and repairs them. Its name comes from the three techniques it
combines, which are abduction~\cite{CoxPietrzykowski86}, belief
revision~\cite{Nebel91} and conceptual change through
reformation~\cite{BundyMitrovic16}. Abduction adds axioms that explain an
observation, belief revision deletes or weakens axioms, and reformation
changes the language of the theory, for example by splitting a predicate in
two or by changing its arity. ABC is written in Prolog and works on
theories in Datalog~\cite{CeriGottlobTanca90}, where theoremhood is
decidable. It has been applied to defeasible reasoning, virtual bargaining,
root cause analysis in mobile networks and the diagnosis of arithmetic
misconceptions~\cite{Bundy26,BundyPhilalithisLi20,LiEtAl22,Tang16}.

Repairing a theory is close in spirit to the analysis of failed proofs in
inductive theorem proving. Proof critics~\cite{IrelandBundy96} patch a
failed proof attempt by speculating a lemma or a generalization, while ABC
patches the axioms. Theory exploration systems such as
IsaScheme~\cite{MontanoRivas12}, IsaCoSy~\cite{JohanssonDixonBundy11},
QuickSpec~\cite{ClaessenSmallboneHughes10} and Hipster~\cite{Hipster14} work
on the other side of the same loop. Given a theory, they find and prove its
interesting properties, which describe what the theory predicts. A proof
assistant that hosts both tasks could diagnose a theory, repair it and then
explore what the repaired theory says. This paper takes a first step in
that direction for Isabelle/HOL~\cite{NipkowPaulsonWenzel02}. We follow the
presentation of ABC in a recent chapter by Bundy~\cite{Bundy26}, whose
Section~5 defines the faults and the repair operations that we implement.

Our contributions are the following.
\begin{itemize}
  \item An implementation of the fault diagnosis and of the ten repair
        operations of ABC as an extension of IsaScheme, in a separate
        Isabelle session that does not modify IsaScheme
        (Section~\ref{sec:extension}). Theories are Horn clauses written in
        Isabelle's syntax and may contain function symbols, so proofs are
        searched up to a depth limit, as Bundy suggests for logics beyond
        Datalog.
  \item A breadth-first search for sequences of repair operations, which
        reports the fault-free theories found with the smallest number of
        operations and with each larger number up to a bound
        (Section~\ref{sec:search}).
  \item The reproduction of the worked examples of the chapter, where the
        repairs described by Bundy appear among the fault-free theories
        with one or two operations, together with the size of the search
        and its cost (Section~\ref{sec:examples}). A closer look at the
        arithmetic example shows that the repair that forgets the carry
        leaves a variable unbound, so the repaired theory predicts any tens
        digit.
  \item A pilot that carries the arithmetic example over to recursive
        functions in Isabelle/HOL. The repair removes an argument of a
        function, and IsaScheme's inductive prover proves what the repaired
        function computes and why the misconception is hard to notice
        (Section~\ref{sec:functions}).
\end{itemize}

Section~\ref{sec:abc} recalls the ABC system, Section~\ref{sec:extension}
describes the extension, Section~\ref{sec:examples} the examples of the
chapter and Section~\ref{sec:functions} the pilot on recursive functions.
Section~\ref{sec:related} discusses related work, and
Section~\ref{sec:conclusion} the limitations of the extension and future
work.

%% ===========================================================================
\section{The ABC system}
\label{sec:abc}

\subsection{Theories and preferred structures}

ABC represents a theory as a finite set of Horn clauses in Kowalski form,
\[
  Q_1 \wedge \dots \wedge Q_m \imp R,
\]
where the $Q_i$ and $R$ are atoms, $m \geq 0$ and the head $R$ may be
missing. A clause with $m = 0$ is an \emph{assertion}, or fact, written
$\imp R$. A clause with $m \geq 1$ and a head is a \emph{rule}, and a clause
without head is a \emph{goal clause}. As in Prolog, names that start with an
upper-case letter are variables, which are implicitly universally
quantified, and the other names are constants and predicates. Datalog has
no function symbols, so a theory has finitely many ground atoms and
theoremhood is decidable.

The observations are given by a \emph{preferred structure} $\mathit{PS}$,
which consists of two sets of ground atoms. The true set $\TPS$ contains
the atoms that should be theorems, and the false set $\FPS$ the atoms that
should not. ABC adopts the closed world assumption. In particular, two
syntactically distinct constants are taken to be unequal, and since
Datalog has no negation, $\neq$ is a predicate of its own. Axioms whose
head is an equation or a disequation can override this assumption, for
instance when an entity has two names.

ABC proves goals by selected literal (SL) resolution~\cite{KowalskiKuehner71}.
One atom of a goal clause is selected and resolved with the head of a rule
or of an assertion, and the process is repeated until the empty goal
clause is reached. Since all the clauses are Horn clauses, each step of a
proof uses an axiom of the theory. A repair can therefore be applied
directly to the axiom involved in a step, without propagating it through
derived clauses.

\subsection{Faults and repairs}

\begin{definition}[Faults]
Let $T$ be a theory and $\mathit{PS}$ a preferred structure. An
\emph{incompatibility} is an atom $\varphi \in \FPS$ such that
$T \vdash \varphi$. An \emph{insufficiency} is an atom $\varphi \in \TPS$
such that $T \nvdash \varphi$. A theory without incompatibilities and
insufficiencies is \emph{fault-free}.
\end{definition}

An incompatibility is repaired by blocking each of its proofs, which
requires making at least one step of each proof fail. An insufficiency is
repaired by letting one or more failed steps of a partial proof succeed.
ABC has five operations for each kind of fault, named after the area of artificial intelligence they come
from. They are summarized in Table~\ref{tab:operations}. For an
incompatibility, the operations target a step that resolves a goal
$P(s_1,\dots,s_n)$ with the head $P(t_1,\dots,t_n)$ of an axiom. For an
insufficiency, they target a failed step between a goal
$P(s_1,\dots,s_m)$ and the head $P'(t_1,\dots,t_n)$ of an axiom, where
either $P \neq P'$ or some $s_i$ and $t_i$ do not unify.

\begin{table}[t]
\centering
\small
\begin{tabularx}{\textwidth}{@{}l X@{}}
\toprule
\multicolumn{2}{@{}l}{\emph{Incompatibility}, goal $P(\bar s)$ resolved
  with the head $P(\bar t)$ of the targeted axiom} \\
\midrule
Belief Revision 1 & Delete the targeted axiom. \\
Belief Revision 2 & Add an unprovable precondition to the body of the
  targeted axiom. \\
Reformation 1 & Rename $P$ to a new predicate $P'$ in the targeted axiom. \\
Reformation 2 & Add an argument to every occurrence of $P$, so that the new
  arguments of the goal and of the targeted head do not unify. \\
Reformation 3 & If $s_i$ and $t_i$ are the same constant $C$, change $t_i$
  to a new constant $C'$. \\
\midrule
\multicolumn{2}{@{}l}{\emph{Insufficiency}, goal $P(\bar s)$ failing
  against the head $P'(\bar t)$ of an axiom} \\
\midrule
Abduction 1 & Add an axiom whose head unifies with the goal. \\
Abduction 2 & Delete from its axiom the body atom that created the goal. \\
Reformation 4 & Replace $P'(\bar t)$ by $P(\bar s)$ in the axiom. \\
Reformation 5 & Remove from every occurrence of $P'$ an argument $i$ such
  that $s_i$ and $t_i$ do not unify. \\
Reformation 6 & If $s_i$ and $t_i$ are distinct constants $C$ and $C'$,
  rename $C'$ to $C$ in all the axioms, or replace that occurrence of $C'$
  by a new variable. \\
\bottomrule
\end{tabularx}
\caption{The repair operations of ABC, after Definitions~5.2 and~5.3 of
Bundy's chapter~\cite{Bundy26}.}
\label{tab:operations}
\end{table}

Reformation~1, 2 and~3 introduce new symbols, which ABC calls dummy terms.
Giving them meaningful names is a symbol grounding problem, for which Chan,
Li and Bundy use large language models~\cite{ChanLiBundy23}. A fault may
need several operations, and many sequences of operations may repair a
theory, so ABC produces several repaired theories. Urbonas et al.\ prune
them to a Pareto front with a Max-SAT solver~\cite{Urbonas20}. For logics
beyond Datalog, Bundy notes that the search for proofs needs a heuristic
limit. An atom of the true set may then be misclassified as an
insufficiency because its proof is too long, and an incompatibility may go
unnoticed for the same reason~\cite{Bundy26}. Wong, Li and Bundy describe
a version of the repair operations for first-order logic~\cite{WongLiBundy24}.

%% ===========================================================================
\section{An extension of IsaScheme}
\label{sec:extension}

The extension lives in the directory \isa{papers/abc-theory-repair} of the
IsaScheme repository. It is an Isabelle session, \isa{ABC\_Repair}, built
on top of IsaLibs, the Isabelle library of IsaScheme, and it does not
change IsaScheme. Its core is a Standard ML file of about 600 lines,
\isa{abc.ML}, loaded by the theory \isa{ABC.thy}, which also declares two
diagnostic commands, \isa{abc\_faults} and \isa{abc\_repair}. The ML
structure \isa{ABC} exports the same functions, so that theories can also
be built programmatically, as we do for the arithmetic example.

\subsection{Language}
\label{sec:language}

Clauses and atoms are written as Isabelle propositions inside quotes. The
conjunction $\wedge$ separates the atoms of a body, and either the
meta-implication $\imp$ or the implication $\longrightarrow$ of HOL
separates the body from the head. A clause whose head is $\mathit{False}$
is a goal clause. The text is parsed with Isabelle's syntax but not
type-checked, so names such as \isa{sum} or \isa{add}, which are constants
of HOL, can be used as predicates without clashes. After parsing, an
unapplied free name that starts with an upper-case letter becomes a
variable, every other name becomes a constant, a predicate or a function
symbol, numerals become constants, and the equality and disequality of HOL
become the built-in predicates $=$ and $\neq$. Internally, a term is a
variable or a symbol applied to a list of terms,
\begin{lstlisting}
datatype fterm = V of indexname | F of string * fterm list
type clause = {body : fterm list, head : fterm option}
\end{lstlisting}
so predicates, function symbols and constants share one representation, and
a predicate is identified by its name and its arity.

The swan theory of Bundy's Example~6.1, with its preferred structure, is
written as follows:
\begin{lstlisting}
abc_faults
  axioms "german X ⟹ euro X" "euro X ∧ swan X ⟹ white X"
         "german bruce" "swan bruce"
  true "black bruce"
  false "white bruce"
\end{lstlisting}
The command prints the faults in the notation of the chapter, with one
proof of each incompatibility and the failed goals of each insufficiency:
\begin{lstlisting}[language=Out]
Faults:
  Insufficiency: black(bruce) is in the true set but is not a theorem.
    Failed goals: black(bruce)
  Incompatibility: white(bruce) is in the false set but is a theorem.
    Proof:
      white(bruce)  by  euro(X) ∧ swan(X) ⟹ white(X)
      euro(bruce)  by  german(X) ⟹ euro(X)
      german(bruce)  by  ⟹ german(bruce)
      swan(bruce)  by  ⟹ swan(bruce)
\end{lstlisting}

\subsection{Proof search}
\label{sec:proofs}

Function symbols make Horn logic undecidable, so the extension cannot
decide theoremhood as ABC does for Datalog. It searches SL resolution
proofs depth first, always selecting the leftmost atom of the goal clause
and trying the clauses in the order of the theory. A branch is cut after
\isa{abc\_max\_depth} resolution steps, 30 by default, and the whole search
stops after \isa{abc\_max\_proofs} proofs, 5 by default. Clauses are
renamed apart with a fresh index at each step, and unification is ordinary
first-order unification with occurs check.

Every step of a proof records the selected goal, the clause it was resolved
with and the origin of the goal, which is the clause and the position in
its body of the atom that created it. A goal that unifies with no head is a
\emph{dead end}, and it is recorded with its origin as well. Branches cut
by the depth limit are not dead ends. These records are what the repair
operations need. A step tells an incompatibility repair which axiom to
target, and a dead end tells an insufficiency repair which goal failed and
which body atom created it.

The built-in predicates are solved without axioms. A goal $s = t$ succeeds
by unifying $s$ and $t$, and a goal $s \neq t$ succeeds when $s$ and $t$
are ground and distinct, which is the closed world assumption of ABC.
Clauses with $=$ or $\neq$ in their heads are used like any other clause,
so they can override the assumption as in ABC.

The function \isa{ABC.faults} searches proofs for each atom of the
preferred structure. An atom of the true set without proofs is an
insufficiency, which keeps the dead ends of its search, and an atom of the
false set with proofs is an incompatibility, which keeps the proofs found.
Insufficiencies are listed before incompatibilities, in the order of the
preferred structure.

\subsection{Repair operations}
\label{sec:operations}

The operations follow Definitions~5.2 and~5.3 of Bundy's chapter. This
section describes the choices that the definitions leave open.

\paragraph{Incompatibilities}
An incompatibility is repaired on the first of its proofs. Each step of the
proof that uses a clause $k$ of the theory yields the following candidates.
Belief Revision~1 deletes clause $k$. Belief Revision~2 adds to its body the
atom $\mathit{dum}(X_1,\dots,X_r)$, where $\mathit{dum}$ is a new predicate
and $X_1,\dots,X_r$ are the variables of the head, so the precondition is
unprovable and the clause can no longer be used. Reformation~1 renames the
predicate of the head to a primed name, as in $\mathit{euro}'$.
Reformation~3 is applied to every argument position where the goal and the
head carry the same constant $C$, which is changed to $C'$ in the head.
Steps solved by the built-in predicates are not targeted, and Reformation~1
and~2 are never applied to $=$ and $\neq$, because renaming the equality or
changing its arity would take its meaning away.

Reformation~2 needs more care, since it changes every occurrence of the
predicate $P$ and must still block the proof. The targeted head gets a new
constant as its extra argument, and the other facts get distinct new
constants too. The clause whose body atom created the goal gets a new
constant when that atom is its only occurrence of $P$. When the clause has
other occurrences of $P$, all of them share a new variable instead. The
targeted step then succeeds, but the shared variable forces the other
occurrences to agree with the new constant of the targeted fact, which
blocks a later step of the proof. This is the repair of the unique mothers
example of the chapter, shown in Section~\ref{sec:mothers}. The remaining
rules get a fresh variable shared by all their occurrences of $P$, so they
keep working for every value of the new argument.

\paragraph{Insufficiencies}
An insufficiency is repaired at its dead ends, of which the first ten
distinct ones are considered. For a dead end with goal $P(s_1,\dots,s_m)$,
Abduction~1 adds the goal as a new fact, whose variables become ordinary
clause variables. Abduction~2 deletes from its clause the body atom that
created the goal, unless the goal is the observation itself. The other
three operations need a \emph{near miss}, a head that almost resolves with
the goal. The chapter does not say how to choose it, and the extension
considers two kinds. A head with the same predicate and arity as the goal
is a near miss when no other head of that predicate fails on fewer
arguments. A head $Q(t_1,\dots,t_m)$ with another predicate of the same
arity is a near miss when its arguments unify with those of the goal.

Reformation~5 removes from every occurrence of $P$ an argument on which a
near miss of the first kind fails. Reformation~6 takes such an argument
when both terms are constants $c$ and $c'$, and either renames $c'$ to $c$
in all the axioms or replaces $c'$ by a new variable in the near miss.
Reformation~4 takes a near miss of the second kind and either replaces $Q$
by $P$ in its head, keeping the arguments $t_i$, which unify with the
$s_i$, or merges $P$ into $Q$ in all the axioms. The second form extends
the definition of the chapter. It is the one used in the root cause
analysis of Section~\ref{sec:5g}, where the predicate $\mathit{ms}$ of a
rule body is a rival name for the predicate $\mathit{microservice}$ of the
facts. The built-in predicates are excluded from Reformation~4 and~5.

New predicates and constants are called $\mathit{dum}_1$,
$\mathit{dum}_2$, and so on, and renamed symbols get a prime. As in ABC,
the preferred structure is never modified, so the new symbols cannot
appear in it.

\subsection{Search for repairs}
\label{sec:search}

A fault may need several operations, and an operation may repair one fault
and create another, so the extension searches sequences of operations
breadth first. Each node of the search is a theory. A node without faults
is a solution and is not repaired further. Otherwise its first fault is
repaired with every applicable operation, and the theories obtained form
the next level, after removing those already seen, which are compared by
their printed clauses. The search stops after \isa{abc\_max\_repairs}
levels, 2 by default. It reports, for each level, the fault-free theories
found with that number of operations, showing at most
\isa{abc\_max\_solutions} of them, 10 by default, with the clauses that
each repair deletes and adds.

Repairing only the first fault of each node fixes the order in which faults
are treated, so the search does not generate the permutations of a
sequence of independent repairs. The remaining faults, including the other
proofs of a repaired incompatibility, are found again when the next level
is diagnosed.

The report shows every level and not only the first one with solutions,
because the smallest repairs are often coarse, such as the deletion of a
rule, while the repair that a human would choose may need one more
operation. The black swan and the 5G examples below are of this kind.
Table~\ref{tab:commands} summarizes the commands and the options, which are
set with Isabelle's attribute syntax, as in
\isa{declare [[abc\_max\_repairs = 3]]}.

\begin{table}[t]
\centering
\small
\begin{tabularx}{\textwidth}{@{}l l X@{}}
\toprule
Command or option & Default & Meaning \\
\midrule
\isa{abc\_faults} & & Print the faults, with a proof of each
  incompatibility and the failed goals of each insufficiency. \\
\isa{abc\_repair} & & Print the faults and the fault-free theories found
  for each number of operations. \\
\isa{abc\_max\_depth} & 30 & Resolution steps of a proof. \\
\isa{abc\_max\_proofs} & 5 & Proofs kept for each atom. \\
\isa{abc\_max\_repairs} & 2 & Operations of a repair. \\
\isa{abc\_max\_solutions} & 10 & Theories shown for each number of
  operations. \\
\bottomrule
\end{tabularx}
\caption{Commands and options of the extension. Both commands take the
arguments \isa{axioms}~$c_1 \dots c_k$ \isa{true}~$a_1 \dots a_l$
\isa{false}~$b_1 \dots b_p$, where the $c_i$ are clauses and the $a_i$ and
$b_i$ atoms, and the last two parts are optional.}
\label{tab:commands}
\end{table}

%% ===========================================================================
\section{The examples of the chapter}
\label{sec:examples}

The theory \isa{ABC\_Examples.thy} contains the worked examples of Bundy's
chapter. Table~\ref{tab:examples} summarizes the search for each of them.
The candidates of a level are the distinct theories generated at that
level, and the fault-free ones are reported as repairs. The last column
gives the level and the position in the report of the repair described in
the chapter. The times are those of \isa{abc\_repair} with the default
options, measured with Isabelle2026 on an Intel Xeon E5645 at 2.40\,GHz.
The search runs in a single thread.

\begin{table}[t]
\centering
\footnotesize
\setlength{\tabcolsep}{4pt}
\begin{tabular}{@{}l r c r r r r r c@{}}
\toprule
& & Faults & \multicolumn{2}{c}{1 operation} & \multicolumn{2}{c}{2 operations}
  & & Chapter \\
\cmidrule(lr){4-5}\cmidrule(lr){6-7}
Example & Clauses & ins./inc. & cand. & free & cand. & free & Time (s)
  & level, pos. \\
\midrule
Black swan (6.1)  &   4 & 1/1 &  5 &  4 &  17 &  17 & $<0.01$ & 2, 7th \\
Mothers (7)       &   3 & 0/1 & 12 & 12 & --- & --- & $<0.01$ & 1, 6th \\
Penguin (9.1)     &   3 & 0/1 & 13 &  8 &  23 &   1 & $<0.01$ & 1, 8th \\
5G, stage 1 (9.3) &   9 & 2/0 & 14 &  3 & 117 &  32 & 16.2    & 2, 6th \\
5G, stage 2 (9.3) &  10 & 0/3 & 14 &  9 &  70 &  70 & $<0.01$ & 1, 5th \\
Arithmetic (9.4)  & 101 & 1/0 & 51 &  6 & 366 & 166 & 0.07    & 1, 5th \\
\bottomrule
\end{tabular}
\caption{The search on the examples of the chapter, whose section is given
in parentheses. The faults are counted as insufficiencies and
incompatibilities. In the unique mothers example every candidate of the first
level is fault-free, so there is no second level.}
\label{tab:examples}
\end{table}

\subsection{The black swan}

Bruce is a German swan, the theory concludes that it is white, and it is
observed to be black. With one operation, four theories are fault-free.
Each applies Reformation~4 to one clause and turns its head into
$\mathit{black}$, either in a rule, as in
$\mathit{german}(X) \imp \mathit{black}(X)$, or in one of the two facts.
Each of them makes $\mathit{black}(\mathit{bruce})$ a theorem and blocks
the proof of $\mathit{white}(\mathit{bruce})$ at the same time. They are
small but implausible, since they change what the theory says about all
Germans or all European swans, or replace a fact about Bruce. The repair of the
chapter needs two operations:
\begin{lstlisting}[language=Out]
Abduction 1: add ⟹ black(bruce)
Reformation 1: rename euro to euro' in german(X) ⟹ euro(X)
  - german(X) ⟹ euro(X)
  + german(X) ⟹ euro'(X)
  + ⟹ black(bruce)
\end{lstlisting}
Here $\mathit{euro}'$ means a European resident and $\mathit{euro}$ a swan
of European type. The repair is the seventh of the 17 repairs with two
operations, all of which start with Abduction~1 because the insufficiency
is the first fault. The other ones block the proof of
$\mathit{white}(\mathit{bruce})$ at another step or with another operation.

\subsection{Unique mothers}
\label{sec:mothers}

The third axiom of the theory
\begin{lstlisting}
abc_repair
  axioms "mum diana william" "mum camilla william"
         "mum X Z ∧ mum Y Z ⟹ X = Y"
  false "diana = camilla"
\end{lstlisting}
allows one mother per person, so the two mothers of William prove that
Diana and Camilla are the same person. The proof resolves the goal with the
rule and then the two body atoms with the two facts. Twelve repairs with one
operation block one of these steps, and the sixth one is Reformation~2 on
the first fact:
\begin{lstlisting}[language=Out]
Reformation 2: increase the arity of mum by one, blocking
               mum(diana, Z_3) against ⟹ mum(diana, william)
  - mum(X, Z) ∧ mum(Y, Z) ⟹ X = Y
  + mum(X, Z, Dum) ∧ mum(Y, Z, Dum) ⟹ X = Y
  - ⟹ mum(diana, william)
  - ⟹ mum(camilla, william)
  + ⟹ mum(diana, william, dum1)
  + ⟹ mum(camilla, william, dum2)
\end{lstlisting}
This is the theory of kinds of mothers of the chapter, before symbol
grounding names $\mathit{dum}_1$ and $\mathit{dum}_2$ \emph{birth} and
\emph{step}. Since the rule has two occurrences of $\mathit{mum}$, its new
argument is a shared variable, so the rule only identifies two mothers of
the same kind. Without the exclusion of the built-in predicates,
Reformation~1 and~2 would also apply to the equality in the head of the
rule and produce repairs without meaning.

\subsection{Defeasible reasoning}

The penguin theory states that birds fly, that penguins are birds and that
Tweety is a penguin, and $\mathit{fly}(\mathit{tweety})$ is in the false
set. The atom $\mathit{penguin}(\mathit{tweety})$ is also put in the true
set, so that the repairs that delete it are rejected. Of the 13 theories
generated with one operation, eight are fault-free, and the last of them
is Reformation~2 on $\mathit{bird}$:
\begin{lstlisting}[language=Out]
  - bird(X) ⟹ fly(X)
  - penguin(X) ⟹ bird(X)
  + bird(X, dum1) ⟹ fly(X)
  + penguin(X) ⟹ bird(X, dum2)
\end{lstlisting}
The new constants play the role of \emph{normal} and \emph{abnormal} in
the chapter, so only normal birds fly and penguins are abnormal birds. The
repaired theory is still monotonic and still says that Tweety is a bird,
which is how ABC models defeasible reasoning without a non-monotonic logic.

\subsection{Root cause analysis of a 5G network}
\label{sec:5g}

The root cause analysis of a 5G network uses ABC in two
stages~\cite{LiEtAl22,Bundy26}. In the first one, the three observed
failures $\mathit{fail}(\mathit{id}_1)$, $\mathit{fail}(\mathit{id}_2)$ and
$\mathit{fail}(\mathit{id}_3)$ are in the true set, but only
$\mathit{fail}(\mathit{id}_2)$ is a theorem. The extension reports two
insufficiencies. Their failed goals include
$\mathit{createOn}(\mathit{id}_1, d_1)$, because no fact says that
$\mathit{id}_1$ is created on $d_1$, and $\mathit{ms}(X, s_1)$, because the
rule for $\mathit{depend}$ uses the predicate $\mathit{ms}$ while the facts
use $\mathit{microservice}$. Three repairs with one operation are
fault-free. One deletes the atom $\mathit{createOn}(Y, X)$ from the rule for
$\mathit{fail}$, one removes the first argument of $\mathit{createOn}$, and
one replaces $\mathit{id}_2$ by a variable in
$\mathit{createOn}(\mathit{id}_2, d_1)$. Each of them makes every service
fail, so they explain the observations without saying why. The repair of
the chapter is the sixth of the 32 repairs with two operations:
\begin{lstlisting}[language=Out]
Abduction 1: add ⟹ createOn(id1, d1)
Reformation 4: merge ms into microservice in all axioms
  - ms(X, s1) ∧ ms(Y, s2) ∧ sameRoute(X, Y) ⟹ depend(Y, X)
  + microservice(X, s1) ∧ microservice(Y, s2) ∧ sameRoute(X, Y)
      ⟹ depend(Y, X)
  + ⟹ createOn(id1, d1)
\end{lstlisting}

Almost all of the 16 seconds of this example are spent on one candidate
theory of the second level, while every other candidate is diagnosed in
less than 0.2 seconds. In that candidate, Abduction~1
adds the dead end $\mathit{ms}(X_{35}, s_1)$ as a fact with a variable, and
Reformation~4 then merges $\mathit{ms}$ into $\mathit{depend}$. The rule
for $\mathit{depend}$ becomes recursive, the new fact makes
$\mathit{depend}(X, s_1)$ true for every $X$, and every attempt to prove a
failure explores the two recursive rules down to the depth limit.

In the second stage, the abduced fact and the merged rule are part of the
theory, and the three failures move to the false set, so each of them is
an incompatibility. Nine theories with one operation are fault-free. The
fifth one deletes the fact $\imp \mathit{full}(d_1)$, which blocks the three
proofs at once and identifies the root cause found in the chapter, namely
that the device $d_1$ is full. The other ones weaken or delete the rule
that derives the failures, or change $\mathit{full}$ or
$\mathit{createOn}$. In the chapter, the second stage is not allowed to
undo the repairs of the first one. The extension has no such protection,
so the repairs with two operations include some that delete the abduced
fact $\mathit{createOn}(\mathit{id}_1, d_1)$.

\subsection{A student's arithmetic misconception}
\label{sec:arith}

Bundy represents the addition of two two-digit numbers as a relation
$\mathit{sum}$ with seven arguments, the two digits of each operand
followed by the three digits of the result, and the addition of two digits
as a relation $\mathit{add}(a, b, u, c)$, where $u$ and $c$ are the units
digit and the carry of $a + b$. The rule
\[
\begin{array}{l}
  \mathit{add}(U_1, U_2, U_3, C_1) \wedge \mathit{add}(T_2, C_1, T_2', 0)
  \wedge \mathit{add}(T_1, T_2', T_3, C_2) \\
  \qquad \imp \mathit{sum}(T_1, U_1, T_2, U_2, C_2, T_3, U_3)
\end{array}
\]
adds the units, adds the carry to the tens digit of the second operand and
then adds the tens. A student who forgets to carry answers $49 + 83 = 22$,
so $\mathit{sum}(4, 9, 8, 3, 0, 2, 2)$ is in the true set. The chapter uses
three digit facts and abduces the missing ones after the repair, and notes
that the whole table of digit facts could be part of the theory. We take
the whole table, generated with the ML interface of the extension:
\begin{lstlisting}
val add_facts : ABC.clause list =
  map_product (fn a => fn b =>
      {body = [], head = SOME (ABC.F ("add", [digit a, digit b,
         digit ((a + b) mod 10), digit ((a + b) div 10)]))})
    (0 upto 9) (0 upto 9)
val sum_rule = ABC.read_clause @{context}
  "add U1 U2 U3 C1 ∧ add T2 C1 T2' 0 ∧ add T1 T2' T3 C2
     ⟹ sum T1 U1 T2 U2 C2 T3 U3"
val _ = writeln (ABC.repair_report @{context}
                  (sum_rule :: add_facts)
                  (preferred ["sum 4 9 8 3 0 2 2"] []))
\end{lstlisting}
The function \isa{digit} builds the constant for a digit, and
\isa{preferred} reads the atoms of the true and the false sets.

The proof fails at the goal $\mathit{add}(4, 9, 2, 0)$, the tens column
once the carry has been added to the 8. Six theories with one operation are
fault-free. Reformation~5 on the fourth argument of $\mathit{add}$, the
carry, is the diagnosis of the chapter, and it needs no abduction here,
because the facts $\mathit{add}(8, 0, 8)$ and $\mathit{add}(4, 8, 2)$ are
already in the table. The other five are Abduction~1, which adds the wrong
fact $\mathit{add}(4, 9, 2, 0)$ and would correspond to a careless slip,
Abduction~2, which deletes the atom for the tens column, Reformation~5 on
the first or the third argument, and Reformation~6, which renames the digit
1 to 0 in all the axioms and so also sets every carry to 0.

A closer look at the theory repaired by Reformation~5 on the carry shows
that it explains too much. After the repair, the rule reads
\[
\begin{array}{l}
  \mathit{add}(U_1, U_2, U_3) \wedge \mathit{add}(T_2, C_1, T_2') \wedge
  \mathit{add}(T_1, T_2', T_3) \\
  \qquad \imp \mathit{sum}(T_1, U_1, T_2, U_2, C_2, T_3, U_3).
\end{array}
\]
The variable $C_1$ is no longer bound by the units column, so it ranges
over all the digits, and $C_2$ occurs only in the head. The repaired theory
therefore proves $\mathit{sum}(4, 9, 8, 3, 0, d, 2)$ for every digit $d$,
and also $\mathit{sum}(4, 9, 8, 3, 7, 2, 2)$, although it still determines
the units digit. The fragment of three facts used in the chapter hides this
problem, because the only fact abduced for the second atom is
$\mathit{add}(8, 0, 8)$.

More observations narrow the diagnosis. A second answer of the student,
$18 + 25 = 33$, leaves four repairs with one operation, since Abduction~1
and Reformation~6 no longer explain both answers. Putting in the false set
the answers 32 and 43, which the student did not give, rejects the
repair by Reformation~5, and no repair is found with at most two
operations. The missing step instantiates the variables left unbound by
Reformation~5 with the carry 0, which gives the rule
\[
\begin{array}{l}
  \mathit{add}(U_1, U_2, U_3) \wedge \mathit{add}(T_2, 0, T_2') \wedge
  \mathit{add}(T_1, T_2', T_3) \\
  \qquad \imp \mathit{sum}(T_1, U_1, T_2, U_2, 0, T_3, U_3).
\end{array}
\]
The extension confirms that this theory has no faults with respect to the
two answers and the two rejected ones. ABC has no operation for this step,
which we also perform by hand on recursive functions in the next section.

%% ===========================================================================
\section{Repairing recursive functions}
\label{sec:functions}

Datalog has no functions, which forces the seven-argument representation
of the previous section. In higher-order logic, the natural representation
of the same knowledge is a recursive function. The theory
\isa{ABC\_Arithmetic.thy} represents numbers as lists of decimal digits,
least significant digit first, and defines column-wise addition with a
carry:
\begin{lstlisting}
fun val :: "nat list ⇒ nat" where
  "val [] = 0"
| "val (d # ds) = d + 10 * val ds"

abbreviation units :: "nat ⇒ nat" where
  "units s ≡ (if s < 10 then s else s - 10)"
abbreviation carry :: "nat ⇒ nat" where
  "carry s ≡ (if s < 10 then 0 else 1)"

fun add_carry :: "nat ⇒ nat list ⇒ nat list ⇒ nat list" where
  "add_carry c [] [] = (if c = 0 then [] else [c])"
| "add_carry c (a # as) [] =
     units (a + c) # add_carry (carry (a + c)) as []"
| "add_carry c [] (b # bs) =
     units (b + c) # add_carry (carry (b + c)) [] bs"
| "add_carry c (a # as) (b # bs) =
     units (a + b + c) # add_carry (carry (a + b + c)) as bs"
\end{lstlisting}
A predicate \isa{valid} states that all the digits are below 10. The
correct theory is right, and it predicts $49 + 83 = 132$:
\begin{lstlisting}
lemma add_carry_correct:
  "c ≤ 1 ⟹ valid xs ⟹ valid ys ⟹
   val (add_carry c xs ys) = c + val xs + val ys"
  by induct_auto
lemma correct_prediction: "add_carry 0 [9, 4] [3, 8] = [2, 3, 1]"
  by simp
\end{lstlisting}
The method \isa{induct\_auto} of IsaLibs searches, with an A* search, for
an induction followed by \isa{auto}, choosing by itself the induction
variable, the induction rule and the variables to generalize. The observation $49 + 83 = 22$ becomes the equation
\isa{add\_carry 0 [9, 4] [3, 8] = [2, 2]}, whose negation the simplifier
proves, so it is an insufficiency of the correct theory. For ground
observations about terminating functions, evaluation decides theoremhood,
and no depth limit is needed.

Reformation~5 removes the carry argument of \isa{add\_carry}. Every call
\isa{add\_carry (carry s) xs ys} becomes a call of the new function, and the
occurrences of \isa{c} that no equation binds any more are instantiated
with 0, which is the step that the Datalog version lacks. Simplifying
$a + b + 0$ to $a + b$ gives the student's procedure, which writes the
units digit of each column and forgets the carry:
\begin{lstlisting}
fun add_student :: "nat list ⇒ nat list ⇒ nat list" where
  "add_student [] [] = []"
| "add_student (a # as) [] = units a # add_student as []"
| "add_student [] (b # bs) = units b # add_student [] bs"
| "add_student (a # as) (b # bs) =
     units (a + b) # add_student as bs"
\end{lstlisting}
The repaired theory explains the observation, since
\isa{add\_student [9, 4] [3, 8] = [2, 2]} is proved by simplification.

At this point theory exploration takes over. A diagnosis is more useful
when we know what the misconceived procedure computes, and IsaScheme
proves it by induction. With the function \isa{carries xs ys}, which lists
the carry of each column, the student's answer loses ten times the value of
the carries:
\begin{lstlisting}
theorem add_student_value:
  "valid xs ⟹ valid ys ⟹
   val (add_student xs ys) + 10 * val (carries xs ys) =
   val xs + val ys"
  by induct_auto
lemma add_carry_differs [simp]:
  "0 < c ⟹ c < 10 ⟹ add_student xs ys ≠ add_carry c xs ys"
  by induct_auto
theorem add_student_correct_iff_no_carry:
  "add_student xs ys = add_carry 0 xs ys ⟷
   val (carries xs ys) = 0"
  by induct_auto
lemma add_student_comm: "add_student xs ys = add_student ys xs"
  by induct_auto
lemma add_carry_comm: "add_carry c xs ys = add_carry c ys xs"
  by induct_auto
\end{lstlisting}
For $49 + 83$ the carries are $[1, 1]$ and $132 - 110 = 22$. The third
theorem explains why the misconception is hard to notice, as the chapter
observes. The student's answer is right exactly when no column produces a
carry. Its inductive proof uses as a simplification rule the second
lemma, which says that a carry changes the next column. The
last two lemmas show that forgetting the carry does not break
commutativity, so the misconception cannot be detected by swapping the
operands. The proofs by \isa{induct\_auto} take about 90, 40, 20, 45, 1 and
15 seconds on 8 threads, in the order of the listings. For the first four,
the theory also gives a proof by the induction or case analysis rule of the
function followed by \isa{auto}.

In this pilot the repair is written by hand. Automating it requires
lifting the operations of ABC from clauses to function definitions.
Reformation~5 on a function removes an argument from every equation and
every call and instantiates the variables that are no longer bound, after
which the function package~\cite{Krauss06} must prove termination again.
Reformation~4 and~6 correspond to replacing a function symbol or a
constant in an equation, and Belief Revision~1 to deleting an equation.
The diagnosis can rely on evaluation for ground observations, and on
IsaScheme's counterexample search with Quickcheck and
Nitpick~\cite{Bulwahn12,BlanchetteNipkow10} for observations with
variables.

%% ===========================================================================
\section{Related work}
\label{sec:related}

ABC grew out of work on the evolution of representations. Bundy and McNeill
argued that representations should be fluents that change with the
reasoning task~\cite{BundyMcNeill06}, and the ORS system of McNeill and
Bundy repaired the ontologies of planning agents after failed plan
executions~\cite{McNeillBundy07}. Reformation, the algorithm behind the
conceptual changes of ABC, was introduced by Bundy and Mitrovic as a
domain-independent way of blocking and unblocking
unifications~\cite{BundyMitrovic16}. Li, Bundy and Smaill combined it with
abduction and belief revision in ABC~\cite{LiBundySmaill18,Li21}. Later
work prunes the repairs with Max-SAT~\cite{Urbonas20}, names the dummy
terms with large language models~\cite{ChanLiBundy23}, extends the
operations to first-order logic~\cite{WongLiBundy24} and applies ABC to
virtual bargaining, root cause analysis and arithmetic
misconceptions~\cite{BundyPhilalithisLi20,LiEtAl22,Tang16}. Our extension
follows the presentation of Bundy's recent chapter~\cite{Bundy26}. We are
not aware of a previous implementation of the operations of ABC inside a
proof assistant.

Belief revision in the AGM tradition~\cite{AGM85} studies the rational
contraction and revision of deductively closed belief sets, and
syntax-based approaches revise finite bases of axioms~\cite{Nebel91}, as
Belief Revision~1 and~2 do. Abduction computes hypotheses that explain an
observation~\cite{CoxPietrzykowski86}, and abductive logic programming
integrates it with Horn clause reasoning~\cite{KakasKowalskiToni92}.
Neither changes the language of the theory, which is the role of
reformation. Lakatos analysed how mathematical concepts are stretched or
split in the face of counterexamples~\cite{Lakatos76}, an idea close to
Reformation~1 and~2.

Theory revision in inductive logic programming refines a Horn clause
domain theory from positive and negative examples~\cite{MuggletonDeRaedt94}.
FORTE~\cite{RichardsMooney95}, for instance, revises a theory by adding and
deleting clauses and antecedents, which correspond to abduction and to
belief revision in ABC, and its examples play the role of the true and
false sets. Algorithmic debugging~\cite{Shapiro83} locates the faulty
clause of a logic program by asking an oracle about intermediate goals,
and model-based diagnosis~\cite{Reiter87} computes minimal sets of faulty
components, which is close to the root cause analysis of
Section~\ref{sec:5g}. The BUGGY models of Brown and Burton~\cite{BrownBurton78}
describe arithmetic misconceptions as perturbations of correct procedures,
which is the view taken by the arithmetic examples of
Sections~\ref{sec:arith} and~\ref{sec:functions}.

In inductive theorem proving, proof critics~\cite{IrelandBundy96} analyse
failed proof attempts to speculate lemmas and generalizations, and
IsaScheme uses this idea in its inductive prover~\cite{MontanoRivas12}.
ABC applies the same productive use of failure to the axioms of a theory.
Theory exploration
systems~\cite{MontanoRivas12,JohanssonDixonBundy11,ClaessenSmallboneHughes10,Hipster14}
discover the properties of a theory, and Section~\ref{sec:functions} uses
IsaScheme to describe the theory produced by a repair.

%% ===========================================================================
\section{Limitations, future work and conclusion}
\label{sec:conclusion}

\paragraph{Proofs outside the kernel}
The extension reasons about Horn theories in ML, outside the logic, so its
proofs are not checked by Isabelle's kernel and its theories are not
Isabelle theories. A Horn theory has a natural representation as an
inductive definition~\cite{Paulson94} with one introduction rule per
clause. The proofs found by SL resolution could be replayed as theorems
about it, and ground atoms could be decided by the predicate
compiler~\cite{BerghoferBulwahnHaftmann09}. Each repaired theory would then
be an Isabelle theory that IsaScheme can explore.

\paragraph{Approximate diagnosis}
With recursion and function symbols, the depth limit makes the diagnosis
approximate. An atom of the true set whose proofs are longer than the limit
is reported as an insufficiency, and an atom of the false set with only
long proofs hides an incompatibility. The 5G example also shows the cost
of the limit, since one candidate with two recursive rules takes almost all
the time of the search. Tabled resolution, as used by Datalog engines,
would make the diagnosis exact again for theories without function symbols
and would avoid the repeated exploration of recursive goals.

\paragraph{Too many repairs}
The search returns many fault-free theories, 166 with two operations in
the arithmetic example, in the order in which they are generated. Larger
theories need a ranking, for example by the Pareto front of Urbonas et
al.~\cite{Urbonas20}, by the plausibility of each operation, or by
protecting the axioms that the user trusts, as the second stage of the root
cause analysis requires.

\paragraph{A missing operation}
The arithmetic example shows that Reformation~5 may leave variables
unbound and produce a theory that predicts too much. An operation that
instantiates such variables with a constant would complete the diagnosis of
the chapter. On recursive functions we applied it by hand, together with
the rest of the repair, and automating the repair of function definitions
along the lines of Section~\ref{sec:functions} is the most interesting
direction for future work, since it would connect the repair of a theory
with its exploration by IsaScheme.

\paragraph{Conclusion}
We have implemented the diagnosis of faults and the ten repair operations
of ABC as an extension of IsaScheme for Isabelle/HOL, for Horn theories
that may contain function symbols. A breadth-first search over sequences of
operations finds the repairs of the worked examples of Bundy's chapter
among the fault-free theories with one or two operations, in less than a
second in all the examples but one. The examples also show the need for
ranking the repairs and for an operation that binds the variables left
free by the removal of an argument. In the pilot on recursive functions,
the repair of the arithmetic misconception becomes a change of a function
definition, after which IsaScheme proves what the student's procedure
computes. Joining theory repair and theory exploration in one proof
assistant is the goal of this line of work.

%% ===========================================================================
\section*{Availability}
{\sloppy
IsaScheme/IsaLibs is available at
\url{https://github.com/omarmrivas/IsaScheme}. The extension and the
theories of Sections~\ref{sec:examples} and~\ref{sec:functions} are in the
directory \url{papers/abc-theory-repair} of the repository. From the
root of the repository, the command
\begin{center}
\isa{isabelle build -d IsaLibs -d papers/abc-theory-repair ABC\_Repair}
\end{center}
builds them in about two minutes on 8 threads.\par}

\bibliographystyle{elsarticle-num}
\bibliography{references}

\end{document}
